\appendix
\renewcommand{\thefigure}{S\arabic{figure}}
\renewcommand{\theHfigure}{S\arabic{figure}}
\setcounter{figure}{0}
\renewcommand{\thetable}{S\arabic{table}}
\renewcommand{\theHtable}{S\arabic{table}}
\setcounter{table}{0}
\section{Data and parameters}
\label{app:data}

%% The input-data inventory (Table 1) appears in the body (Section~\ref{sec:data}).
%% This appendix holds the algorithm-parameter table and the reproducibility
%% checklist for the analysis pipeline.

\subsection{Parameters}
\label{app:params}
Table~\ref{tab:params} lists the settings used in the analysis. All are logged with
the released code so that the chain is reproducible end to end.

\begin{table}[htbp]
  \centering\footnotesize
  \begin{singlespace}
  \caption{Algorithm parameters and settings (full-census pipeline).}\label{tab:params}
  \begin{tabular}{@{}p{0.34\textwidth}p{0.58\textwidth}@{}}
  \toprule
  Stage / parameter & Value \\
  \midrule
  Coordinate reference system & EPSG:32635 (WGS 84 / UTM zone 35N, central meridian 27$^\circ$E, false easting 500{,}000\,m) \\
  Analysis grid             & 250\,m square fishnet (INSPIRE/GHSL-aligned); 500\,m for the scale test \\
  Cells intersecting region & 16{,}496 \\
  Urban-cell exclusion      & built fraction $\ge 0.10$; water fraction $\le 0.50$; mean slope $\le 15\%$ \\
  Urban cells retained      & 3{,}777 \\
  Land-cover fractions      & ESA WorldCover 2021 (10\,m), via \texttt{exactextract} \\
  Distance covariates       & to coastline (OSM); to core (Konak, 27.1287\,$^\circ$E, 38.4192\,$^\circ$N) \\
  Fabric strata             & 7 a-priori, rule-based (provisional) \\
  Analysis set              & full census of all retained urban cells ($N=3{,}777$; no sampling) \\
  Building / street source  & official \.{I}zmir IMM layers (901{,}609 buildings; 207{,}734 segments) \\
  Floor-area ratio          & from measured floor counts (field \texttt{ZEMINUSTUK}) \\
  Network service areas     & 400\,m and 800\,m network-distance reaches along the IMM road graph (capped Dijkstra from the nearest node) \\
  Orientation entropy       & 18 bins, length-weighted \\
  Land-surface temperature  & Landsat JJA mean 2014--2024, 100\,m, \si{\celsius} (zonal mean per cell) \\
  Feature scaling           & $z$-score after winsorizing to the [1, 99] percentiles \\
  Dimension reduction       & PCA, smallest set with $\ge 90\%$ variance (7 components, 91\% realized) \\
  Clustering                & Ward linkage; $k$ by silhouette subject to minimum cluster size $\ge 10 \Rightarrow k=4$ \\
  Scale stability           & re-cluster at 500\,m; Adjusted Rand Index vs 250\,m ($=0.38$; 250\,m $k=4$, 500\,m $k=3$) \\
  Heat model                & gradient boosting (XGBoost: 400 trees, depth 4, learning rate 0.05, subsample 0.8) \\
  Cross-validation          & spatial-block: 6 $k$-means blocks on $(x,y)$, GroupKFold (6 folds) \\
  Explanation               & TreeSHAP; global and per-stratum mean $|\text{SHAP}|$ \\
  Social vulnerability      & $z$(elderly share) $+\ z$(age dependency), ADNKS 2024; 3{,}777/3{,}777 cells matched (centroid-in-\emph{mahalle}; no fallback) \\
  Spatial weights           & $k$-nearest neighbours ($k=8$), row-standardized \\
  Spatial statistics        & Moran's $I$, LISA, Getis-Ord $G_i^{*}$ (999 permutations, seed 0) \\
  Need axes (priority)      & heat, cooling deficit, access deficit, coastal exposure, social vulnerability \\
  Aggregation               & TOPSIS, equal and entropy weights \\
  Weight robustness         & Monte-Carlo, 2{,}000 Dirichlet$(\mathbf{1})$ weightings \\
  Cell-level frontier       & Pareto non-dominance over the 3{,}777 cells (223 retained) \\
  \bottomrule
  \end{tabular}
  \end{singlespace}
\end{table}

\subsection{Reproducibility}
\label{app:repro}
\begin{itemize}
  \item \textbf{Code.} The complete analysis code, including data acquisition and
  conversion, is openly released in a version-controlled repository.
  \item \textbf{Environment.} Python 3.12 with pinned dependencies (GeoPandas, OSMnx,
  momepy, libpysal/esda, scikit-learn, XGBoost, SHAP, pymoo, earthengine-api).
  \item \textbf{Determinism.} Fixed seeds throughout, covering clustering, the SHAP
  spatial blocks and the spatial-statistics permutations (seed 0), and a seeded
  Monte-Carlo generator for the weight robustness.
  \item \textbf{Data provenance.} All inputs are open or openly licensed and are
  catalogued with their sources, resolutions and licences in
  Table~\ref{tab:data-inventory}.
  \item \textbf{Parameter logging.} All run parameters are retained with the released
  outputs.
  \item \textbf{Provisional classification.} The rule-based seven-stratum
  stratification is open to refinement with planning input; the unsupervised
  clustering tests rather than assumes it.
\end{itemize}

\section{Supplementary analyses}
\label{app:supp_analyses}

\subsection{Aggregation dependence}
\label{app:cluster-priority}
To test whether the equity inversion is a property of the planning-oriented stratum partition or a
scale-invariant property of urban form, we re-ran the full priority pipeline (five need axes, entropy and
Monte-Carlo TOPSIS) on the four unsupervised morphometric super-types instead of the seven strata
(Table~\ref{tab:cluster-priority}). The cluster ordering is statistically orthogonal to the stratum ordering:
correlating the two priority ranks at the cell level (each cell inheriting its stratum rank and its cluster rank)
gives Spearman $\rho = -0.06$ ($p < 0.001$). The waterfront signal does not survive aggregation to clusters
because the 249 waterfront cells scatter 42\%/33\%/24\% across three clusters with a plurality in none; at the
cluster level the entropy weighting collapses onto access deficit (0.64) and social vulnerability (0.32) with
negligible heat weight, and the coolest, most vulnerable coastal fabric is no longer resolvable as a unit. This
confirms that the inversion is a decision-frame-dependent finding tied to the strata in which planning operates,
and should not be read as invariant to the spatial unit of analysis.

\begin{table}[htbp]
  \centering\footnotesize
  \caption{Adaptation priority computed on the four unsupervised morphometric super-types (full census, $N=3{,}777$),
  for comparison with the stratum-level ranking of Table~\ref{tab:priority}. Columns as in Table~\ref{tab:priority};
  rows ordered by Monte-Carlo mean rank. The cluster order is near-orthogonal to the stratum order (cell-level
  Spearman $\rho = -0.06$, $p < 0.001$).}
  \label{tab:cluster-priority}
  \begin{tabular}{@{}lrrrrr@{}}
    \toprule
    Morphometric super-type & $n$ & TOPSIS$_\text{eq}$ & TOPSIS$_\text{ent}$ & Mean rank & Rank [p10--p90] \\
    \midrule
    Sparse open edge              & 123   & 0.75 & 0.99 & 1.3 & 1--3 \\
    Compact mid-rise (dense core) & 1{,}265 & 0.36 & 0.18 & 1.8 & 1--2 \\
    Coarse-grain large-footprint  & 1{,}062 & 0.17 & 0.02 & 3.2 & 2--4 \\
    Low-rise open periphery       & 1{,}327 & 0.11 & 0.08 & 3.7 & 3--4 \\
    \bottomrule
  \end{tabular}
\end{table}

\subsection{Heat leverage}
\label{app:heat-leverage}
The main-text priority ranking uses measured summer LST as the heat axis. Because part of a fabric's thermal
signature is attributable to fixed geography (distance to coast, distance to core, slope) rather than to
modifiable urban form, we decomposed each stratum's mean SHAP heat attribution into a \emph{modifiable} component
(built form, floor-area ratio, coverage, greenery and street/network descriptors; 12 features) and a \emph{fixed}
component (distance to coast, distance to core and slope; 3 features), and re-ran the priority pipeline substituting
the modifiable-heat leverage for measured LST on the heat axis (Table~\ref{tab:heat-leverage},
Figure~\ref{fig:supp-heat-leverage}). The equity inversion is robust to this substitution: waterfront remains the
top priority (Monte-Carlo mean rank $1.60 \rightarrow 1.95$), historic core second, peripheral expansion third,
and the overall stratum ordering is preserved (Spearman $\rho = 0.86$, $p = 0.014$). The one material change is
that industrial/logistics rises from fourth to approximately second: it is the only stratum whose summer heat is
driven by modifiable built form ($+1.18\,\si{\celsius}$ modifiable SHAP), whereas waterfront's coolness is almost
entirely fixed geography ($-1.64\,\si{\celsius}$ coastal SHAP). This decomposition sharpens, rather than
overturns, the planning reading: waterfront needs adaptation on vulnerability and exposure grounds, while
industrial fabric offers the greatest heat leverage through form-based intervention.

\begin{table}[htbp]
  \centering\scriptsize
  \setlength{\tabcolsep}{3pt}
  \caption{Modifiable-versus-fixed decomposition of each stratum's mean SHAP heat attribution ($\si{\celsius}$),
  and the resulting priority-rank shift when modifiable-heat leverage replaces measured LST on the heat axis.
  Modifiable = built form, FAR, coverage, greenery, street/network (12 features); Fixed = distance to coast,
  distance to core, slope (3 features). Heat leverage is the positive part of the modifiable component. The
  overall ordering is preserved (Spearman $\rho = 0.86$, $p = 0.014$).}
  \label{tab:heat-leverage}
  \begin{tabular}{@{}lrrrrrr@{}}
    \toprule
    Fabric stratum & Modifiable & Fixed & Heat leverage & Raw LST & MC rank raw & MC rank leverage \\
    \midrule
    Waterfront transformation & $-0.09$ & $-1.64$ & 0.00 & 41.68 & 1.60 & 1.95 \\
    Historic core             & $-0.23$ & $-0.39$ & 0.00 & 42.74 & 2.71 & 2.99 \\
    Peripheral expansion      & $-0.52$ & $-0.53$ & 0.00 & 42.38 & 2.98 & 3.42 \\
    Industrial / logistics    & $1.18$  & $0.71$  & 1.18 & 45.46 & 4.15 & 2.46 \\
    Apartment block           & $-0.06$ & $0.31$  & 0.00 & 43.66 & 5.10 & 5.37 \\
    Hillside / incremental    & $-0.32$ & $-0.03$ & 0.00 & 43.02 & 5.68 & 5.92 \\
    Grid residential          & $-0.09$ & $0.00$  & 0.00 & 43.24 & 5.78 & 5.89 \\
    \bottomrule
  \end{tabular}
\end{table}

\subsection{Vulnerability robustness}
\label{app:vuln-robust}
The main-text social-vulnerability axis is a demographic (age-structure) composite, and might be criticized as a
narrow operationalization of vulnerability. To test whether the equity inversion is an artefact of that choice, we
derived two additional, age-independent vulnerability axes from open Global Human Settlement Layer (GHSL, JRC
R2023A, 100\,m) products and re-ran the full priority pipeline under each. (i) A \emph{building-age} proxy: the
area-weighted mean construction epoch per cell, obtained from the GHS-BUILT-S multi-temporal built-up-surface
stack (between 1975 and 2020); older stock indexes both a poorer thermal envelope and
pre-modern-code construction, which is relevant to the wider seismic-risk programme. (ii) A \emph{dwelling-crowding} proxy: GHS-POP 2020 residents per unit of
estimated residential floor area (building footprint area $\times$ count $\times$ mean floors, from our own
morphometrics). Both axes are near-orthogonal to the age-based index at the cell level (Spearman $\rho = +0.13$
building age, $\rho = -0.19$ crowding), confirming they carry independent information rather than restating age
structure.

The equity inversion is robust to the vulnerability definition (Table~\ref{tab:vuln-robust},
Figure~\ref{fig:supp-vuln-robust}). Under the building-age axis the waterfront remains the top-ranked stratum
(Monte-Carlo mean rank $2.04$) and the full ordering is nearly identical to the age-based baseline (Spearman
$\rho = 0.96$). Under a composite that averages all three axes (age, building age, crowding) the waterfront is again
first ($1.52$). The one definition under which the waterfront is displaced is dwelling crowding, where it falls to a
near-tie for second (rank $2.88$) behind peripheral expansion ($2.82$). This is intelligible because crowding
peaks in the densely-occupied peripheral growth fabric rather than on the amenity-driven waterfront, and one that
identifies crowded peripheral expansion as the fabric whose priority is most sensitive to a crowding-based reading
of need. In no definition does the cool, high-amenity waterfront leave the top tier; the paradox that the coolest
fabric carries the highest adaptation priority therefore does not depend on an age-based construction of
vulnerability.

\begin{table}[htbp]
  \centering\footnotesize
  \caption{Monte-Carlo mean priority rank of each fabric stratum under four social-vulnerability definitions:
  the main-text age-structure composite, a GHSL building-age proxy, a GHS-POP dwelling-crowding proxy, and a
  three-axis composite. Rows ordered by the main (age-SVI) ranking. Waterfront stays first under three of four
  definitions and second under crowding; the building-age ranking reproduces the age-based ranking almost exactly
  (Spearman $\rho = 0.96$).}
  \label{tab:vuln-robust}
  \begin{tabular}{@{}lrrrr@{}}
    \toprule
    Fabric stratum & Age-SVI (main) & Building age & Crowding & Combined \\
    \midrule
    Waterfront transformation & 1.59 & 2.04 & 2.88 & 1.52 \\
    Historic core             & 2.70 & 2.33 & 4.79 & 2.71 \\
    Peripheral expansion      & 2.98 & 4.04 & 2.82 & 3.50 \\
    Industrial / logistics    & 4.16 & 3.92 & 2.88 & 3.85 \\
    Apartment block           & 5.07 & 4.12 & 4.91 & 4.60 \\
    Hillside / incremental    & 5.70 & 5.54 & 5.84 & 6.31 \\
    Grid residential          & 5.79 & 6.01 & 3.89 & 5.50 \\
    \bottomrule
  \end{tabular}
\end{table}

\subsection{SHAP stability}
\label{app:shap-stability}
Because the explainable heat model is a screening tool with modest out-of-sample skill (spatial-block
cross-validated $R^2 = 0.27$), we verified that the SHAP attributions on which the mechanism interpretation rests
are stable across space rather than an artefact of any single region. Refitting the model on each of the six
spatial-block training folds and recomputing SHAP values on the held-out block, the feature-importance ordering is
highly stable (mean pairwise Spearman $= 0.86$ across folds; minimum $0.77$; Figure~\ref{fig:supp-shap-stability}a).
More important for interpretation, every high-importance predictor retains the same sign of effect in all six folds: distance to
coast and greenery are cooling, building-footprint area and slope are warming, in every fold
(Figure~\ref{fig:supp-shap-stability}b). The only predictors whose sign is inconsistent across folds,
building coverage and intersection density, are among the least important and carry none of the reported conclusions. One
spatial block generalizes poorly (held-out $R^2 = -0.88$; median fold $R^2 \approx 0.22$), which is expected of a
screening model and is the reason we do not read the model as predictive; the attributions that support the
planning interpretation are nonetheless spatially robust.

\section{Supplementary figures}
\label{app:supp_figs}

\begin{figure}[htbp]
  \centering
  \includegraphics[width=\textwidth]{supp_m1_heat_leverage}
  \caption{Modifiable-versus-fixed heat decomposition and its effect on priority ranking.
  (a) Per-stratum mean SHAP contribution to summer LST, split into a modifiable component (built form and greenery)
  and a fixed component (coast distance, core distance, slope); only industrial/logistics has positive modifiable
  heat leverage, whereas waterfront's coolness is almost entirely fixed coastal geography ($-1.64\,\si{\celsius}$).
  (b) Monte-Carlo mean priority rank under the measured-LST heat axis versus the SHAP-coupled modifiable-heat
  leverage axis; waterfront stays first and the overall ordering is preserved (Spearman $\rho = 0.86$), with
  industrial/logistics the only stratum to shift materially (fourth to $\approx$ second).}
  \label{fig:supp-heat-leverage}
\end{figure}

\begin{figure}[htbp]
  \centering
  \includegraphics[width=\textwidth]{supp_vulnerability_robustness}
  \caption{Robustness of the adaptation-priority ranking to the social-vulnerability
  definition. (a) Monte-Carlo mean priority rank of each fabric stratum across four vulnerability definitions
  (age-structure composite, GHSL building age, GHS-POP dwelling crowding, and a three-axis composite); the
  waterfront (bold) stays top-ranked under three of four and second under crowding. (b) Cell-level Spearman
  correlations showing the two GHSL-derived axes are near-orthogonal to the age-based index ($|\rho| < 0.2$),
  i.e.\ they encode independent dimensions of vulnerability rather than restating age structure.}
  \label{fig:supp-vuln-robust}
\end{figure}

\begin{figure}[htbp]
  \centering
  \includegraphics[width=\textwidth]{supp_shap_stability}
  \caption{Spatial stability of the SHAP heat attributions across six spatial-block
  cross-validation folds. (a) Per-fold mean $|$SHAP$|$ for each predictor (blue dots, one per held-out block; red
  bar = mean); the importance ordering is fold-stable (rank Spearman $= 0.86$). (b) Directional sign-consistency:
  fraction of folds in which each predictor's effect keeps the same sign. The high-importance predictors (coast
  distance, greenery, footprint area, slope; starred) are directionally identical in all six folds, whereas the
  sign-unstable predictors are all low-importance.}
  \label{fig:supp-shap-stability}
\end{figure}

\begin{figure}[htbp]
  \centering
  \includegraphics[width=\textwidth]{supp_data_quality}
  \caption{Spatial resolution and input data characteristics. (a) Sensor point spread function (PSF) and grid footprint commensurability between the Landsat thermal product and INSPIRE grids. (b) Boxplot of official IMM building footprint counts across the urban census by fabric stratum. (c) Bivariate geographic distribution of sample cells by distance to the urban core and distance to the coastline, colored by stratum. (d) Bivariate demographic input covariation showing the relationship between elderly (65+) population share and the age dependency ratio across the sample, colored by stratum.}
  \label{fig:supp-data-quality}
\end{figure}

\begin{figure}[htbp]
  \centering
  \includegraphics[width=\textwidth]{supp_pca_diagnostics}
  \caption{PCA dimensionality reduction diagnostics. (a) Scree plot showing individual and cumulative explained variance ratios across principal components, with a dashed line at the 90\% threshold. (b) PCA component loadings biplot for PC1 and PC2, showing the directional alignment of feature vectors. (c) Feature loadings heatmap for PC1 to PC7, illustrating which morphological and environmental variables load onto each component. (d) Grid sample cells plotted in principal component space (PC1 vs. PC2), color-coded by their unsupervised fabric-resilience cluster membership ($k=4$) to show clustering separation.}
  \label{fig:supp-pca-diagnostics}
\end{figure}

\begin{figure}[htbp]
  \centering
  \includegraphics[width=\textwidth]{supp_xgb_diagnostics}
  \caption{Explainable heat model XGBoost regression diagnostics. (a) Cross-validation learning curves showing RMSE for training and validation sets across boosting rounds (up to 400 trees), averaged over the 6 spatial-block folds. (b) Density curves of out-of-fold prediction residuals (Actual LST $-$ Predicted LST) comparing the Full Model and the Morphology-Only Model. (c) Observed vs. out-of-fold predicted summer LST scatter plot for the Full Model, color-coded by fabric stratum and overlaid with a 1:1 reference line. (d) Spatial map of prediction residuals (Actual $-$ Predicted LST) for the Full Model plotted at grid cell coordinates, color-coded by error magnitude to diagnose spatial patterns.}
  \label{fig:supp-xgb-diagnostics}
\end{figure}

\begin{figure}[htbp]
  \centering
  \includegraphics[width=\textwidth]{supp_topsis_sensitivity}
  \caption{Multi-objective TOPSIS ranking weight sensitivity. (a) Dirichlet distribution of weight fractions across the five adaptation axes in the Monte Carlo simulation ($2{,}000$ runs). (b) Kendall's rank correlation matrix showing cell-level rank agreement between equal, entropy, heat-dominant, access-dominant and social-dominant weighting regimes. (c) Rank sensitivity of the peripheral-expansion stratum as a function of the weight allocated to the access-deficit axis. The redundant strata closeness-distribution panel was removed because main Figure~\ref{fig:topsis-robustness}a reports the corresponding rank distributions.}
  \label{fig:supp-topsis-sensitivity}
\end{figure}
